\chapter{Patterns: derived state}
\label{ch:derived}

\begin{aside}
Many UI values are derived from others. Total of a list,
filtered subset, formatted display string. Memos are the
tool.
\end{aside}

\section{Computed values}

\begin{lstlisting}
auto items = signal<std::vector<Item>>(...);
auto total = memo([&] {
    int sum = 0;
    for (auto& i : ctx.get(items))
        sum += i.value;
    return sum;
});
\end{lstlisting}

The memo re-runs only when items changes. Other code reads
the memo as a signal.

\section{Chained memos}

\begin{lstlisting}
auto items = signal<...>(...);
auto query = signal<std::string>("");
auto filtered = memo([&] {
    return filter(ctx.get(items), ctx.get(query));
});
auto count = memo([&] {
    return ctx.get(filtered).size();
});
\end{lstlisting}

A change to query triggers filtered, then count. Effects on
count fire if and only if count actually changed.

\section{When not to memo}

Cheap computations don't need memoisation. If the function
is ``return x + 1'', just inline.

The memo machinery costs some bytes per entry; use when
the computation is non-trivial or the subscribers are many.

\section{Async derived state}

For derivations that need I/O, use an async effect that
writes to a signal:

\begin{lstlisting}
ctx.async_effect([&]() -> Task<void> {
    auto id = ctx.get(user_id);
    auto data = co_await fetch(id);
    ctx.set(user_data, data);
});
\end{lstlisting}

\section{Recap}

\begin{recap}
\begin{itemize}[leftmargin=*]
  \item Memo for sync derivations.
  \item Async effect for I/O-bound derivations.
  \item Chain memos to layer derivations.
  \item Don't memo trivia.
\end{itemize}
\end{recap}

\section*{Workshop}

\begin{enumerate}[leftmargin=*]
  \item Identify three derived values in your app.
  \item Convert each to a memo.
  \item Profile before and after.
\end{enumerate}
